\section{Ai đo được chỗ thắt}
\label{sec:ch06-ai-do}

\index{chỗ thắt!ai đo được}%
Chương~\ref{ch:encoder-decoder} kết thúc bằng một chuyện chưa xong. Bản kể quen
thuộc nói encoder-decoder hỏng ở câu dài, và chương ấy chỉ ra rằng bài của
Sutskever và cộng sự nói ngược lại ở năm chỗ khác nhau. Nhưng \tn{chỗ
thắt}{bottleneck} thì có thật, chương ấy đo được nó, nên câu hỏi còn lại là ai
đo được nó \emph{trước}, trên một mô hình thật, ở quy mô thật.

Câu trả lời là Cho và cộng sự, trong bài SSST-8 năm 2014 mà
mục~\ref{sec:ch05-cau-noi} gọi là bài 2014b~\autocite{cho2014properties}. Chương
trước mới đọc tóm tắt của nó và nói thẳng ra là chỉ đọc tóm tắt. Phiên này đọc
hết, và phần thân bài đắt hơn tóm tắt nhiều.

\subsection*{Một cái bẫy công cụ, kể ra vì suýt hỏng cả mục này}

\index{đọc bài báo!số trang công cụ báo sai}%
Tôi tải hai bản, bản arXiv:1409.1259v2 và bản của ACL Anthology. Lệnh
\code{file} báo cả hai đều \emph{sáu} trang, và lần trích đầu tiên của tôi cũng
ra sáu trang, đứt giữa câu ở \enquote{In fact, if we limit the lengths of both
the source}. Tôi suýt viết mục này với sáu trang ấy.

Cây trang của hai file bị chẻ làm hai nhánh, và đếm số \code{/Type /Page} trong
byte thô thì ra chín. Ba trang bị bỏ sót chứa hình 4, hình 5, hình 6, toàn bộ
mục 5.2 và phần kết luận. Nói cách khác, đúng phần mà chương này đi tìm.

Bài học không phải về định dạng PDF. Nó là: một công cụ báo cho bạn kích thước
của thứ bạn đang đọc thì con số ấy cũng là một phép đo, và phép đo nào cũng có
thể sai. Cách kiểm rẻ nhất là kiểm bằng nội dung, tức đọc câu cuối cùng của
trang cuối cùng và hỏi nó có phải một câu kết thúc không.

\subsection*{Hai bảng và hai hình}

\index{Cho 2014b!số đo theo độ dài câu}%
Bài huấn luyện hai mô hình trên WMT'14 Anh-Pháp, chỉ lấy cặp câu mà cả hai phía
dài tối đa 30 từ, \tn{từ điển}{vocabulary} 30 nghìn từ thường gặp nhất mỗi phía.
Một mô hình là RNN Encoder-Decoder của Cho 2014a, tức một
\tn{mạng hồi quy}{recurrent neural network}, viết tắt là RNN, ở cả hai đầu; một
là kiến trúc tích chập đệ quy có cổng của chính bài. Đối chứng là Moses, một hệ
dịch thống kê theo cụm từ.

Bảng 1 của bài, điểm Bilingual Evaluation Understudy, viết tắt là BLEU, trên
tập test, cắt lấy phần chương này dùng:

\begin{minted}{text}
                  moi do dai      cau 10-20 tu
RNNenc  toan bo     13.92            20.99
Moses   toan bo     33.30            32.00
RNNenc  khong UNK   23.45            27.03
Moses   khong UNK   35.63            35.40
\end{minted}

Đọc theo hàng thì thấy khoảng cách giữa hai hệ co lại rất mạnh khi giới hạn độ
dài câu: trên toàn bộ tập test, Moses hơn RNNenc 19.38 điểm; trên riêng câu 10
tới 20 từ, còn 11.01. Bỏ luôn các câu chứa từ \tn{ngoài từ điển}{out-of-vocabulary}
thì còn 8.37. Và bài siết thêm một nấc nữa trong phần chữ, ở câu mà tôi cho là
câu đắt nhất của cả bài: giới hạn cả câu nguồn lẫn bản dịch tham chiếu trong
khoảng 10 tới 20 từ \emph{và} bỏ mọi câu có từ ngoài từ điển, thì hai hệ được
27.81 và 33.08.

Hình 4(a) vẽ BLEU của RNNenc theo độ dài câu, làm trơn bằng cửa sổ 10. Đường
lên tới đỉnh khoảng 20 điểm ở quãng câu 10 tới 15 từ, rồi đi xuống, và tới câu
70 tới 80 từ thì gần chạm 0. Hình 5 vẽ đúng phép đo ấy cho Moses. Đường của
Moses đi ngang, hơi lên, nằm trong dải 30 tới 35 suốt mọi độ dài.

\index{Cho 2014b!vì sao hai hình phải đọc cùng nhau}%
Hai hình ấy phải đọc cùng nhau, và đây là chỗ bài này mạnh hơn hẳn một bài chỉ
báo cáo hình 4. Một mình hình 4 chứng minh được rất ít: câu dài thì khó dịch
hơn, chuyện đó ai cũng đoán được, và một đường đi xuống có thể chỉ đang đo độ
khó của bài toán. Hình 5 chặn đường đọc đó lại. Cùng dữ liệu, cùng tập test,
cùng cách chấm điểm, một hệ không dùng mạng nơ-ron: đường của nó không đi xuống.
Vậy cái đi xuống trong hình 4 là tính chất của \emph{kiến trúc}, không phải của
bài toán.

Cả ba hình đều là hình vẽ và không có bảng số nào đứng sau. Nên những con số
tôi vừa đưa ra khi mô tả chúng, đỉnh khoảng 20 và dải 30 tới 35, là số đọc bằng
mắt trên trục và chỉ dùng để nói hình dạng; đừng trích lại chúng như số của bài.
Số của bài trong phần độ dài câu chỉ có hai, 27.81 và 33.08, và chúng là hai
con số duy nhất ở đó được viết thành chữ.

\subsection*{Phản biện hiển nhiên, và câu bài dùng để chặn nó}

\index{Cho 2014b!chặn phản biện về độ dài huấn luyện}%
Mô hình huấn luyện trên câu tối đa 30 từ, rồi bị đem đo trên câu 70 từ. Phản
biện tự đến: nó có bao giờ thấy câu dài thế đâu mà đòi nó dịch được. Nếu đúng
vậy thì hình 4 đo trí nhớ của tập huấn luyện chứ không đo chỗ thắt.

Bài viết một câu ngắn ngay sau đó: \enquote{Note that we observed a similar
trend even when we used sentences of up to 50 words to train these models}.
Huấn luyện lại với câu tới 50 từ, xu hướng vẫn thế. Một câu, không có bảng đi
kèm, và nó đóng lại đúng lối thoát dễ chịu nhất.

Tôi để ý câu này vì đó là kiểu câu chỉ có ở người đã tự đặt câu hỏi ấy cho
mình. Nó không làm bài đẹp hơn và không ai bắt họ chạy thêm lần ấy.

\subsection*{Hai cách đo cùng một thứ}

\index{chỗ thắt!hai cách đo}%
Đặt cạnh nhau thì chương~\ref{ch:encoder-decoder} và bài 2014b đo cùng một
chuyện bằng hai cái núm ngược nhau. Bài 2014b giữ nguyên mô hình và kéo dài câu.
Chương trước giữ nguyên câu và bóp hẹp \tn{vector ngữ cảnh}{context vector}. Cả
hai cách đều là cách thay đổi tỷ lệ giữa lượng thông tin phải đi qua và bề rộng
chỗ nó đi qua, và mục~\ref{sec:ch05-con-lai-gi} đã phát biểu chỗ thắt đúng như
một tỷ lệ chứ không như một tính chất của kiến trúc.

Điều đó cũng nói vì sao hai bài báo không đá nhau. Sutskever và cộng sự chạy
với 8000 số thực cho một câu và không chạm vào chỗ thắt: bốn tầng, mỗi tầng một
nghìn cell, và \tn{ô nhớ}{memory cell} Long Short-Term Memory, viết tắt là
LSTM, mang hai trạng thái. Bài 2014b chạy RNNenc một tầng
một nghìn đơn vị có cổng, mà đơn vị có cổng chỉ mang một trạng thái, nên vector
của nó là một nghìn số. Ít hơn tám lần, và chạm.

\subsection*{Bài Bahdanau phỏng đoán chứ không đo}

\index{Bahdanau 2015!phỏng đoán chứ không đo}%
Chỗ này đáng cẩn thận, vì rất dễ gán cho bài chính của chương một việc mà nó
không làm. Tóm tắt của nó viết:

\begin{quote}
\enquote{we conjecture that the use of a fixed-length vector is a bottleneck in
improving the performance of this basic encoder-decoder architecture}
\end{quote}

\enquote{Conjecture} là chữ của bài. Bài này không thu hẹp vector để xem chuyện
gì xảy ra, cũng không vẽ đường BLEU theo độ dài cho một mô hình vector cố định
của riêng nó ở phần nêu vấn đề. Nó dẫn bằng chứng của người khác. Phần kết
luận, mục 7, ghi rõ hai nguồn: bài 2014b và Pouget-Abadie và cộng
sự~\autocite{pougetabadie2014overcoming}.

Cả hai đều là bài workshop ở SSST-8, cùng nhóm, cùng năm. Nghĩa là toàn bộ nền
thực nghiệm cho câu \enquote{vector cố định là một chỗ thắt}, tại thời điểm bài
2015 được viết, gồm hai bài workshop. Tôi nói ra không phải để hạ chúng xuống,
mà vì cách kể thông thường làm ngược lại: nó gán phép đo cho bài nổi tiếng và
để hai bài workshop chìm đi.

\subsection*{Chỗ bài 2014b tự chỉ sai người}

\index{Cho 2014b!hai câu chỉ vào hai chỗ khác nhau}%
Còn một chuyện nữa trong bài 2014b, và nó chỉ hiện ra khi đọc cả ba trang cuối.

Mục 5.1 chỉ vào vector:

\begin{quote}
\enquote{The most obvious explanatory hypothesis is that the fixed-length
vector representation does not have enough capacity to encode a long sentence
with complicated structure and meaning.}
\end{quote}

Còn kết luận, mục 6, chỉ vào chỗ khác:

\begin{quote}
\enquote{Despite the radical difference in the architecture between RNN and
grConv which were used as an encoder, both models suffer from the curse of
sentence length. This suggests that it may be due to the lack of
representational power in the decoder.}
\end{quote}

Hai encoder rất khác nhau mà hỏng như nhau, nên phần đáng ngờ là phần dùng
chung, tức decoder. Lập luận ấy chặt chẽ và nó chỉ nhầm người. Bài 2015 vá ở
phía encoder: giữ lại mọi trạng thái encoder và cho decoder chọn giữa chúng.
Decoder không hề được nới rộng thêm chút nào.

Chỗ này đáng đọc kỹ hơn là đáng cười. Ba trong bốn tác giả của bài 2014b, Cho,
Bahdanau và Bengio, là tác giả của bài 2015, và khoảng cách giữa hai bài là vài
tháng. Nghĩa là cùng một nhóm người, trong cùng một mùa, viết ra một suy đoán
rồi tìm ra lời giải nằm ở chỗ khác. Suy đoán ấy không phải suy đoán cẩu thả:
nó đúng theo dữ liệu họ có, và cái nó thiếu là một khả năng thứ ba mà chưa ai
nghĩ tới lúc ấy, rằng có thể bỏ hẳn ràng buộc bề rộng cố định đi thay vì chọn
xem nới bên nào.

Mục sau là cách bỏ.
